\section{Bases within a density-zero class}\label{sec:local-bases}

We now construct families of candidate bases within one equivalence class
modulo $\mathcal J_0$. The purpose is twofold: balanced finite changes
must preserve the family, and intervals of subsets must contain or lie
on one side of a member of the family. The latter property will supply
maximal independent extensions in the eventual matroid.

Fix $T\subseteq E_0$ such that $T$ and $T^c$ are
$\mathcal J_0$-positive, and put
\[
 \mathcal C(T)=\{X\subseteq E_0:X\triangle T\in\mathcal J_0\}.
\]
We use $T$ as the chosen representative, or \emph{prototype}, of this class.
Every member of $\mathcal C(T)$ and its complement are
$\mathcal J_0$-positive. Indeed, an ideal-small member, together with its
ideal-small symmetric difference from $T$, would make $T$ ideal-small;
the same argument applies to complements.

For $n\ge1$ and $m\ge0$, define
\begin{align}
 V_n&=\bigcup_{m<n}W_m,\qquad c_n=|V_n|,\notag\\
 p_n^T(X)&=|X\cap V_n|-|T\cap V_n|,\notag\\
 d_m^T(X)&=\frac{|X\cap W_m|-|T\cap W_m|}{|W_m|},\notag\\
 f_n^T(X)&=p_n^T(X)+c_n\left(
       \sup_{m\ge n}d_m^T(X)+\inf_{m\ge n}d_m^T(X)\right).
       \label{eq:probes}
\end{align}
These quantities are defined for $X\in\mathcal C(T)$. The prefix count
records a finite modification exactly once its affected blocks lie in
$V_n$. The tail extrema allow later-block changes to influence earlier
probes, and their sum changes sign when both $X$ and $T$ are complemented.
Write
\[
 \ell_T(X)=\liminf_{n\to\infty}f_n^T(X),\qquad
 h_T(X)=\limsup_{n\to\infty}f_n^T(X),
\]
allowing infinite values, and define two families
\begin{align}
 \mathcal B^-(T)&=\{X\in\mathcal C(T):-1<\ell_T(X)\le0\},\notag\\
 \mathcal B^+(T)&=\{X\in\mathcal C(T):0\le h_T(X)<1\}.
 \label{eq:local-bases}
\end{align}
We call these the lower-limit and upper-limit rules, respectively.
In this section, a basis for a rule means a member of the indicated
family; no matroid is asserted yet.

\subsection{Probe estimates and continuity}

\begin{lemma}[Probe properties]\label{lem:probe-properties}
For $X,Y\in\mathcal C(T)$ the following statements hold.
\begin{enumerate}
\item $d_m^T(X)\to0$ and $p_n^T(X)/c_n\to0$.
\item If $X\subseteq Y$, then
\begin{equation}\label{eq:probe-monotone}
 f_n^T(Y)-f_n^T(X)\ge |(Y\setminus X)\cap V_n|.
\end{equation}
\item Adding or deleting one element changes every probe by at most
one. A finite modification with net change $b$ in cardinality satisfies
$f_n^T(Y)=f_n^T(X)+b$ for all sufficiently large $n$.
\item Complementation gives
\[
 f_n^{T^c}(X^c)=-f_n^T(X),\qquad
 \mathcal B^+(T^c)=\{B^c:B\in\mathcal B^-(T)\}.
\]
\end{enumerate}
\end{lemma}

\begin{proof}
The first density limit follows from $X\triangle T\in\mathcal J_0$.
Also $p_n^T(X)=\sum_{m<n}|W_m|d_m^T(X)$. For any $\varepsilon>0$,
the terms after a fixed index have absolute density at most
$\varepsilon$; their contribution divided by $c_n$ is at most
$\varepsilon$, and the finite initial contribution divided by $c_n$
tends to zero.

Inclusion increases each block count, hence both tail extrema; the
prefix increases by $|(Y\setminus X)\cap V_n|$. This proves the second
assertion. If an element changes in a block $W_k$ with $k<n$, only the
prefix changes. If $k\ge n$, each tail extremum changes by at most
$1/|W_k|$. By \eqref{eq:block-growth},
\[
 |W_k|\ge 2\sum_{m<k}|W_m|\ge 2c_n,
\]
so the total change is at most one. After every block affected by a
finite modification is in the prefix, the tails agree and the prefix
changes by exactly $b$.

Finally, complementing both $X$ and $T$ negates every prefix and block
density. The supremum of the negatives is the negative infimum, and
vice versa. This proves the last assertion, including its half-open
endpoint conventions.
\end{proof}

The next continuity statement is needed both to stop finite deletion
procedures and to pass to their decreasing limit. Its common
density-zero bound is essential to controlling the tail extrema.

\begin{lemma}[Continuity within a fixed interval]\label{lem:block-continuity}
Suppose $I\subseteq X_k\subseteq Y$ for $k\ge1$, where
$I,Y\in\mathcal C(T)$, and membership of every element in $X_k$
eventually agrees with its membership in a set $X$. Then
$X\in\mathcal C(T)$ and, for every fixed $n\ge1$,
\[
 f_n^T(X_k)\longrightarrow f_n^T(X).
\]
\end{lemma}

\begin{proof}
We have $I\subseteq X\subseteq Y$ and $Y\setminus I\in\mathcal J_0$,
so $X\in\mathcal C(T)$. On every fixed finite union of blocks, $X_k$
eventually agrees with $X$. On the remaining blocks,
\[
 |d_m^T(X_k)-d_m^T(X)|
 \le \frac{|(Y\setminus I)\cap W_m|}{|W_m|},
\]
uniformly in $k$, and the right-hand side tends to zero. Thus the
sequences of block densities converge uniformly in $m$. Supremum and
infimum over a fixed tail change by at most the uniform difference,
while the fixed prefix eventually stabilizes. Formula
\eqref{eq:probes} gives the result.
\end{proof}

\begin{lemma}[Antichains and finite changes]\label{lem:local-antichain}
Each family $\mathcal B^-(T)$ and $\mathcal B^+(T)$ contains $T$, is
an antichain under inclusion, and is preserved by finite modifications
that add and delete equally many elements.
\end{lemma}

\begin{proof}
The probes of $T$ vanish. Balanced finite changes preserve the class
and eventually preserve every probe, so they preserve either rule.

Suppose $X\subsetneq Y$ both satisfy the lower-limit rule. A finite
difference of cardinality $k\ge1$ gives
$\ell_T(Y)=\ell_T(X)+k>0$, a contradiction. For an infinite difference,
$|(Y\setminus X)\cap V_n|\to\infty$. The probes of $X$ are eventually
bounded below because their lower limit is finite, so
Lemma~\ref{lem:probe-properties} implies $f_n^T(Y)\to\infty$, again
a contradiction. The upper-limit case follows by complementing and
using the lower-limit rule for $T^c$.
\end{proof}

\subsection{Modifications from a positive pool}

To interpolate between subsets, we need to move the limiting probes
in either direction while remaining in $\mathcal C(T)$. A positive
pool of available elements suffices, even though we use only a
density-zero subset of it.

\begin{lemma}[Small modifications with divergent probes]\label{lem:positive-pool}
Let $X\in\mathcal C(T)$.
If $G\subseteq X^c$ is $\mathcal J_0$-positive, there is
$P\subseteq G$ in $\mathcal J_0$ such that
$f_n^T(X\cup P)\to+\infty$.
If $G\subseteq X$ is $\mathcal J_0$-positive, there is
$P\subseteq G$ in $\mathcal J_0$ such that
$f_n^T(X\setminus P)\to-\infty$.
\end{lemma}

\begin{proof}
For additions, choose $\delta>0$ such that $G$ has density at least
$\delta$ in arbitrarily late blocks. Suppress the superscript $T$,
put $a_n=\sup_{m\ge n}|d_m(X)|$, and choose strictly increasing
positive integers $N_j$ such that, for $n\ge N_j$,
\begin{equation}\label{eq:spike-cutoff}
 \frac{|p_n(X)|}{c_n}+2a_n+\frac1{\sqrt{c_n}}\le 2^{-j}.
\end{equation}
This is possible by Lemma~\ref{lem:probe-properties}.

For each sufficiently large $j$, put $b_j=2^{1-j}<\delta$ and choose
an increasing sequence of blocks $m_j\ge N_{j+1}$ in which $G$ has
density at least $\delta$, also requiring
$1/|W_{m_j}|\le\delta-b_j$. From $G\cap W_{m_j}$ select exactly
$\lceil b_j|W_{m_j}|\rceil$ elements. Rounding is possible because
this number is at most $\delta|W_{m_j}|$. Let $P$ be the union of these
selections. Its nonzero block densities lie between $b_j$ and
$b_j+1/|W_{m_j}|$, and hence tend to zero. Thus $P\in\mathcal J_0$.

If $N_j\le n<N_{j+1}$, the selected block $m_j$ belongs to the tail
in \eqref{eq:probes}. The tail supremum for $X\cup P$ is at least
$b_j-a_n$, and its tail infimum is at least $-a_n$. The prefix can only
increase. Therefore \eqref{eq:spike-cutoff} gives
\[
 \frac{f_n(X\cup P)}{c_n}
 \ge -\frac{|p_n(X)|}{c_n}+b_j-2a_n
 \ge \frac1{\sqrt{c_n}}.
\]
These intervals contain every sufficiently large $n$, proving
divergence to $+\infty$. Apply the addition statement to $X^c$ and
$T^c$, and use complement symmetry, to obtain the deletion statement.
\end{proof}

\subsection{Interpolation and interval alternatives}

The finite-change rule rounds any finite limiting value into the
interval $(-1,0]$. The remaining case has opposite infinite limits at
the two endpoints. We handle it by finite deletions whose effect on
earlier probes is summable.

\begin{lemma}[Interpolation]\label{lem:local-interpolation}
If $I\subseteq Y$ are in $\mathcal C(T)$,
$\ell_T(I)\le-1$, and $\ell_T(Y)>0$, then there is
$B\in\mathcal B^-(T)$ with $I\subseteq B\subseteq Y$.
\end{lemma}

\begin{proof}
Suppress $T$ from the notation. If $\ell(I)=a$ is finite, add
$\lfloor-a\rfloor$ elements from $Y\setminus I$. There are enough:
a smaller finite gap of size $h$ would give
$\ell(Y)=a+h\le-1$. The resulting lower limit belongs to $(-1,0]$.
If $\ell(Y)=b$ is finite, delete $\lceil b\rceil$ gap elements from
$Y$. A smaller finite gap of size $h$ would instead give
$\ell(I)=b-h>0$. Again the resulting lower limit belongs to $(-1,0]$.

It remains to consider $\ell(I)=-\infty$ and $\ell(Y)=+\infty$.
Set $X_0=Y$ and construct decreasing sets $X_j$ by finite deletions.
At stage $j\ge1$, do nothing if some $n\ge j$ already satisfies
$f_n(X_{j-1})\le0$. Otherwise choose $M_j\ge j$ so large that
\begin{equation}\label{eq:stage-loss}
 2c_j\sup_{m\ge M_j}
       \frac{|(Y\setminus I)\cap W_m|}{|W_m|}<2^{-j}.
\end{equation}
This is possible because $Y\setminus I\in\mathcal J_0$.
Delete elements of $X_{j-1}\setminus I$ in blocks $m\ge M_j$, one
at a time in increasing block order, stopping as soon as
$f_n\le0$ for some $n\ge j$.

This stopping time is finite. If every eligible element were deleted,
the resulting set would differ from $I$ by finitely many elements,
so it would have lower limit $-\infty$. Fix a strictly negative probe
with index $n\ge j$ for that set. By
Lemma~\ref{lem:block-continuity}, the same probe is negative after
finitely many of the deletions. Just before the stopping deletion all
probes with index at least $j$ were positive. Each drops by at most
one, so
\begin{equation}\label{eq:stage-reset}
 f_n(X_j)\ge-1\qquad(n\ge j)
\end{equation}
at every stage that makes deletions. For $n<j$ the prefix is unchanged
during the whole stage, and each tail extremum decreases by at most
the density supremum in \eqref{eq:stage-loss}. Since $c_n\le c_j$,
the entire stage decreases $f_n$ by less than $2^{-j}$.

Let $X_\infty=\bigcap_{j\ge0}X_j$. It lies between $I$ and $Y$ and
belongs to $\mathcal C(T)$; Lemma~\ref{lem:block-continuity} applies
to this decreasing sequence. Every stage leaves a witness
$f_n\le0$ with $n\ge j$. Such a witness persists under later
deletions and in the limit, so $\ell(X_\infty)\le0$.

There is a first stage $j_0$ that makes deletions. Otherwise the
unchanged set $Y$ would have a nonpositive probe in every tail,
contrary to $\ell(Y)=+\infty$. Fix $n\ge j_0$. After stage $j_0$ its
probe is at least $-1$. Through stage $n$, every later stage either
leaves it unchanged or restores this bound by \eqref{eq:stage-reset}.
Stages $j>n$ then cost less than $2^{-j}$ each. Passing to the limit,
\[
 f_n(X_\infty)\ge-1-\sum_{j>n}2^{-j}=-1-2^{-n}.
\]
Consequently $-1\le\ell(X_\infty)\le0$. If the lower limit exceeds
$-1$, take $B=X_\infty$. If it equals $-1$, add one point of
$Y\setminus X_\infty$ to obtain lower limit zero. There are infinitely
many such points: a finite difference from $Y$ would contradict
$\ell(Y)=+\infty$ by the finite-change rule. This completes the
interpolation.
\end{proof}

The local rules have the following interval property. Unless a basis
lies below the lower endpoint, between the endpoints, or above the upper
endpoint, the prototype has a positive part outside the lower endpoint
and the upper endpoint has a positive part outside the prototype.
These are the two differences needed to separate a new prototype class
from this one in Section~\ref{sec:assembly}.

\begin{proposition}[Interval alternatives for one class]
\label{prop:local-intervals}
Fix $\sigma\in\{-,+\}$ and let $I\subseteq Y\subseteq E_0$.
If
\[
 T\setminus I\in\mathcal J_0
 \quad\text{or}\quad
 Y\setminus T\in\mathcal J_0,
\]
there is $B\in\mathcal B^\sigma(T)$ satisfying at least one of
\begin{equation}\label{eq:local-alternatives}
 B\subseteq I,\qquad I\subseteq B\subseteq Y,
 \qquad Y\subseteq B.
\end{equation}
\end{proposition}

\begin{proof}
We first establish the interval alternatives when one or both endpoints
belong to the class. For the lower-limit rule, if
$X\in\mathcal C(T)$ has $\ell_T(X)\le-1$,
Lemmas~\ref{lem:positive-pool} and~\ref{lem:local-interpolation}
extend it to a basis by additions from any prescribed positive
subset of $X^c$. If $\ell_T(X)>0$, their deletion versions give a
basis contained in $X$, using any prescribed positive subset of $X$.
Both unrestricted corrections are available because $X$ and $X^c$
are positive.

If both $I$ and $Y$ belong to $\mathcal C(T)$, a lower endpoint with
positive lower limit supplies a basis below $I$, and an upper
endpoint with lower limit at most $-1$ supplies a basis above $Y$.
An endpoint that is itself a basis also gives an outcome in
\eqref{eq:local-alternatives}. The only remaining case is
$\ell_T(I)\le-1$ and $\ell_T(Y)>0$, when
Lemma~\ref{lem:local-interpolation} gives a basis between them.

If $I$ belongs to the class and $Y\setminus I$ is positive,
correcting $I$ downwards, leaving it fixed when it is a basis, or
correcting it upwards from the pool $Y\setminus I$ gives a basis
below $I$ or between $I$ and $Y$.
If $Y$ belongs to the class and the gap is positive, the
corresponding corrections give a basis between $I$ and $Y$ or
above $Y$.
Complementation proves these same three conclusions for the
upper-limit rule, exchanging the two one-endpoint conclusions.

Now suppose $T\setminus I\in\mathcal J_0$.
If $I\in\mathcal C(T)$, a $\mathcal J_0$-small gap puts $Y$ in the same class,
whereas a positive gap permits the lower-endpoint conclusion just
proved. In either case we obtain \eqref{eq:local-alternatives}.
If $I\notin\mathcal C(T)$, then $T\cap I\in\mathcal C(T)$ and
$I\setminus T$ is positive. Applying the lower-endpoint conclusion
to $[T\cap I,I]$ gives a basis contained in $I$.

Finally, if $Y\setminus T\in\mathcal J_0$, apply the result just
proved to the prototype $T^c$ and the interval $[Y^c,I^c]$, with
the opposite rule. Complementing its basis gives
\eqref{eq:local-alternatives} for $T$ and the original rule.
\end{proof}
